\scanpage{120}
\mspage{114}
% archive sheet number: 119

\deleted{Les deux familles \unsure{parties|petites|finies}} $\{A,B\}$ est une paire admissible de parties de $\Sigma_{\unsure{\mathcal A|A}}$.

Ceci n'est pas tout à fait satisfaisant, car on veut pouvoir traiter le cas d'un ``diagramme'' quelconque dans $(\unsure{\mathcal F|\mathscr F},\ll)$
\[
F_0\ll F_{\unsure{2|1}}\ll\cdots\ll F_n\deleted{\unsure{\ll}},
\]
donnant naissance dans $\Sigma=\Sigma_{\mathcal A}$ à un diagramme similaire dans $\unsure{\operatorname{Fig\,spat}|\operatorname{Fig\,opt}}(\Sigma)$, mais \deleted{\unsure{voudrait}}\inserted{qu'on} \unsure{aimerait} bien interpréter, par oubli des structures d'ordre \unsure{superfétatoires} sur \deleted{\unsure{les}} chacun des $F_{\unsure{i|j}}$, comme un diagramme dans $\operatorname{Fig\,spat\,disc}(\Sigma)$
\[
A_0\ll A_1\ll\cdots\ll A_n.
\]

Pour exprimer l'axiome pertinent, \deleted{disons de façon}\inserted{posons de façon} générale la

\uline{Définition} Soit $\Sigma$ un \unsure{spatios|specton}, \deleted{\unsure{Une}}\deleted{\unsure{ens. d'}} $A$ une partie de $\Sigma$. Nous dirons que $A$ est une partie \uline{commensurable} de $\Sigma$, ou que \deleted{\unsure{les a}}
